% ============================================================
\chapter{\'Ethique, Biais et S\'ecurit\'e}
\label{ch:ethique}
% ============================================================

\begin{intuition}[title=\textbf{Responsabilit\'e en TAL}]
Les mod\`eles de langage sont d\'eploy\'es \`a grande \'echelle et
influencent des d\'ecisions cons\'equentes~: recrutement, justice,
sant\'e, \'education. Ils reproduisent et amplifient les biais pr\'esents
dans leurs donn\'ees d'entra\^inement. Comprendre, mesurer et att\'enuer
ces biais est une responsabilit\'e fondamentale du praticien en TAL.
\end{intuition}

% ------------------------------------------------------------
\section{Biais dans les plongements de mots}
\label{sec:ethique:embeddings}
% ------------------------------------------------------------

\begin{definition}[Biais dans les embeddings]
\index{biais!embeddings}
Les plongements de mots capturent les r\'egularit\'es statistiques du
corpus, y compris les st\'er\'eotypes sociaux. Par exemple, dans
Word2Vec entra\^in\'e sur Google News~:
\[
    \vec{\text{homme}} - \vec{\text{femme}} \approx
    \vec{\text{informaticien}} - \vec{\text{m\'enag\`ere}}
\]
\end{definition}

\begin{theoreme}[WEAT --- Word Embedding Association Test]
\index{WEAT}
Le test WEAT (Caliskan et al., 2017) mesure l'association entre deux
ensembles de mots cibles $X, Y$ et deux ensembles d'attributs $A, B$~:
\[
    d(X, Y, A, B) = \sum_{x \in X} s(x, A, B) - \sum_{y \in Y} s(y, A, B)
\]
o\`u $s(w, A, B) = \frac{1}{\abs{A}} \sum_{a \in A} \cos(w, a)
- \frac{1}{\abs{B}} \sum_{b \in B} \cos(w, b)$.
La significativit\'e statistique est \'evalu\'ee par un test de
permutation.
\end{theoreme}

\begin{exemple}[Biais de genre]
En appliquant WEAT avec $X = \{\text{math, science, ing\'enierie}\}$,
$Y = \{\text{arts, litt\'erature, humanit\'es}\}$,
$A = \text{noms masculins}$, $B = \text{noms f\'eminins}$, on observe
une association significative entre sciences et masculin, reproduisant
un st\'er\'eotype de genre.
\end{exemple}

% ------------------------------------------------------------
\section{Biais dans les mod\`eles de langue}
\label{sec:ethique:lm_bias}
% ------------------------------------------------------------

\begin{definition}[Biais des LLM]
\index{biais!LLM}
Les LLM exhibent des biais \`a plusieurs niveaux~:
\begin{enumerate}
    \item \textbf{Repr\'esentationnel}~: certains groupes sont
    sous-repr\'esent\'es ou st\'er\'eotyp\'es.
    \item \textbf{Allocatif}~: le mod\`ele performe mieux pour certains
    groupes (ex.~: r\'esum\'es de CV favorisant certains pr\'enoms).
    \item \textbf{G\'en\'eratif}~: production de contenu toxique,
    discriminatoire ou inexact concernant certains groupes.
\end{enumerate}
\end{definition}

\begin{proposition}[Sources de biais]
\begin{itemize}
    \item \textbf{Donn\'ees d'entra\^inement}~: le web contient des
    contenus biais\'es, sexistes, racistes.
    \item \textbf{Annotation}~: les annotateurs apportent leurs propres
    biais culturels et linguistiques.
    \item \textbf{Objectif d'entra\^inement}~: maximiser la vraisemblance
    reproduit les patterns dominants.
    \item \textbf{D\'eploiement}~: les biais de feedback (les utilisateurs
    interagissent plus avec certains contenus) amplifient les biais
    existants.
\end{itemize}
\end{proposition}

% ------------------------------------------------------------
\section{M\'etriques d'\'equit\'e}
\label{sec:ethique:metriques}
% ------------------------------------------------------------

\begin{definition}[\'Egalit\'e des chances (Equalized Odds)]
\index{equalized odds}
Un classifieur satisfait l'\textbf{\'egalit\'e des chances} si sa
performance est identique pour tous les groupes d\'emographiques $g$~:
\[
    P(\hat{Y} = 1 | Y = y, G = g) = P(\hat{Y} = 1 | Y = y)
    \quad \forall g, \forall y \in \{0, 1\}
\]
\end{definition}

\begin{definition}[Parit\'e d\'emographique]
\index{parit\'e d\'emographique}
La \textbf{parit\'e d\'emographique} exige que la d\'ecision soit
ind\'ependante du groupe~:
\[
    P(\hat{Y} = 1 | G = g) = P(\hat{Y} = 1) \quad \forall g
\]
\end{definition}

\begin{attention}[title=\textbf{Incompatibilit\'e des crit\`eres}]
Il est g\'en\'eralement impossible de satisfaire simultan\'ement
la parit\'e d\'emographique, l'\'egalit\'e des chances et la calibration
(th\'eor\`eme d'impossibilit\'e de Chouldechova, 2017, et Kleinberg
et al., 2016). Le choix du crit\`ere d\'epend du contexte applicatif.
\end{attention}

\begin{definition}[Biais de toxicit\'e]
\index{toxicit\'e}
Le \textbf{biais de toxicit\'e} mesure la propension d'un mod\`ele \`a
g\'en\'erer du contenu toxique (insultes, menaces, discours haineux).
Le score RealToxicityPrompts (Gehman et al., 2020) mesure la probabilit\'e
de g\'en\'erer du contenu toxique \`a partir d'amorces neutres.
\end{definition}

% ------------------------------------------------------------
\section{Techniques de d\'ebiaisage}
\label{sec:ethique:debiaisage}
% ------------------------------------------------------------

\begin{definition}[D\'ebiaisage des embeddings]
\index{d\'ebiaisage}
La m\'ethode de Bolukbasi et al.\ (2016) identifie la direction de
biais $\mathbf{b}$ (ex.~: direction homme-femme) et projette les
embeddings pour \'eliminer cette composante~:
\[
    \mathbf{w}_{\text{debias}} = \mathbf{w} -
    (\mathbf{w} \cdot \mathbf{b}) \mathbf{b}
\]
pour les mots neutres (non d\'efinitionnellement genre\'es).
\end{definition}

\begin{remarque}
Le d\'ebiaisage g\'eom\'etrique des embeddings a \'et\'e critiqu\'e~:
Gonen et Gold (2019) montrent que les biais persistent dans la structure
des voisinages m\^eme apr\`es projection. Des approches plus r\'ecentes
interviennent au niveau des donn\'ees ou de l'entra\^inement.
\end{remarque}

\begin{definition}[Strat\'egies de d\'ebiaisage pour les LLM]
\index{d\'ebiaisage!LLM}
\begin{enumerate}
    \item \textbf{Curation des donn\'ees}~: filtrage des contenus
    toxiques, \'equilibrage des repr\'esentations.
    \item \textbf{Instruction tuning}~: fine-tuning avec des instructions
    explicites d'\'equit\'e.
    \item \textbf{RLHF/DPO}~: entra\^iner le mod\`ele \`a pr\'ef\'erer
    les r\'eponses non biais\'ees via feedback humain.
    \item \textbf{Prompting}~: instructions syst\`eme sp\'ecifiant un
    comportement \'equitable.
    \item \textbf{D\'ecodage contr\^ol\'e}~: modifier les probabilit\'es
    de g\'en\'eration pour r\'eduire la toxicit\'e.
\end{enumerate}
\end{definition}

% ------------------------------------------------------------
\section{D\'etection de toxicit\'e}
\label{sec:ethique:toxicite}
% ------------------------------------------------------------

\begin{definition}[Classifieur de toxicit\'e]
\index{d\'etection de toxicit\'e}
Un \textbf{classifieur de toxicit\'e} est un mod\`ele supervis\'e
entra\^in\'e \`a d\'etecter le contenu offensant, haineux ou dangereux.
Exemples~: Perspective API (Google), OpenAI Moderation API.
\end{definition}

\begin{proposition}[D\'efis de la d\'etection de toxicit\'e]
\begin{enumerate}
    \item \textbf{D\'efinition subjective}~: la toxicit\'e d\'epend du
    contexte, de la culture et de l'intention.
    \item \textbf{Biais du classifieur}~: les mod\`eles de toxicit\'e
    eux-m\^emes peuvent \^etre biais\'es (ex.~: classifier le dialecte
    afro-am\'ericain comme plus toxique).
    \item \textbf{\'Evasion}~: les utilisateurs contournent les filtres
    par des reformulations cr\'eatives.
    \item \textbf{Surtoxicit\'e}~: un filtrage trop agressif censure du
    contenu l\'egitime.
\end{enumerate}
\end{proposition}

% ------------------------------------------------------------
\section{Impact environnemental}
\label{sec:ethique:environnement}
% ------------------------------------------------------------

\begin{definition}[Co\^ut carbone de l'entra\^inement]
\index{impact environnemental}
L'entra\^inement d'un grand mod\`ele de langue consomme une \'energie
consid\'erable. Strubell et al.\ (2019) estiment l'empreinte carbone
de l'entra\^inement d'un Transformer \`a $\sim$284 tonnes de CO$_2$
(comparable \`a 5 fois la vie d'une voiture am\'ericaine).
\end{definition}

\begin{formules}[title=\textbf{Estimation du co\^ut carbone}]
\[
    \text{CO}_2 \approx \text{PUE} \times \text{kWh} \times
    \text{intensit\'e carbone (g CO}_2\text{/kWh)}
\]
o\`u PUE (Power Usage Effectiveness) $\approx 1.1$--$1.4$ pour les
datacenters modernes.
\end{formules}

\begin{proposition}[Strat\'egies de r\'eduction]
\begin{itemize}
    \item \textbf{Mod\`eles plus petits}~: distillation, pruning,
    quantification.
    \item \textbf{Entra\^inement efficient}~: mixed precision, gradient
    checkpointing, architecture search.
    \item \textbf{Localisation}~: entra\^iner dans des r\'egions \`a
    \'electricit\'e d\'ecarbon\'ee.
    \item \textbf{R\'eutilisation}~: partager les mod\`eles pr\'e-entra\^in\'es
    (HuggingFace Hub).
\end{itemize}
\end{proposition}

% ------------------------------------------------------------
\section{Pratiques responsables en IA}
\label{sec:ethique:pratiques}
% ------------------------------------------------------------

\begin{definition}[Model Cards]
\index{model cards}
Une \textbf{model card} (Mitchell et al., 2019) est une fiche
documentant un mod\`ele~: usage pr\'evu, limitations, biais connus,
m\'etriques d'\'evaluation, donn\'ees d'entra\^inement, et
consid\'erations \'ethiques.
\end{definition}

\begin{definition}[Datasheets for Datasets]
\index{datasheets}
Les \textbf{datasheets} (Gebru et al., 2021) documentent les jeux de
donn\'ees~: motivation, composition, processus de collecte,
pr\'etraitement, utilisations recommand\'ees, et distributions
d\'emographiques.
\end{definition}

\begin{remarque}
Les r\'eglementations \'evoluent rapidement~:
\begin{itemize}
    \item \textbf{AI Act} (Union Europ\'eenne, 2024)~: classification
    des syst\`emes d'IA par niveau de risque, obligations de transparence.
    \item \textbf{Executive Order on AI} (\'Etats-Unis, 2023)~: exigences
    de s\'ecurit\'e et de tests pour les mod\`eles de fondation.
    \item Les entreprises adoptent des \emph{Responsible AI frameworks}
    int\'egrant \'evaluation des risques, audits r\'eguliers et comit\'es
    d'\'ethique.
\end{itemize}
\end{remarque}

% ------------------------------------------------------------
\section{Impl\'ementation Python}
\label{sec:ethique:code}
% ------------------------------------------------------------

\begin{codeblock}[title=\textbf{D\'etection de biais dans les embeddings}]
\begin{minted}{python}
import numpy as np
from gensim.models import KeyedVectors

def weat_score(model, X, Y, A, B):
    """Compute WEAT effect size."""
    def s(w, A_set, B_set):
        a_sims = [model.similarity(w, a) for a in A_set
                  if a in model]
        b_sims = [model.similarity(w, b) for b in B_set
                  if b in model]
        return np.mean(a_sims) - np.mean(b_sims)

    x_scores = [s(x, A, B) for x in X if x in model]
    y_scores = [s(y, A, B) for y in Y if y in model]
    d = np.mean(x_scores) - np.mean(y_scores)
    std = np.std(x_scores + y_scores)
    return d / std if std > 0 else 0.0

# Example: gender bias in professions
X = ["programmer", "engineer", "scientist"]
Y = ["nurse", "teacher", "librarian"]
A = ["man", "male", "boy", "he"]
B = ["woman", "female", "girl", "she"]
# score = weat_score(model, X, Y, A, B)
\end{minted}
\end{codeblock}

\begin{codeblock}[title=\textbf{D\'ebiaisage g\'eom\'etrique}]
\begin{minted}{python}
def debias_embedding(embedding, bias_direction):
    """Project out the bias direction from an embedding."""
    projection = np.dot(embedding, bias_direction) * bias_direction
    return embedding - projection

# Compute bias direction (e.g., he - she)
bias_dir = model["he"] - model["she"]
bias_dir = bias_dir / np.linalg.norm(bias_dir)

# Debias a neutral word
debiased = debias_embedding(model["programmer"], bias_dir)
\end{minted}
\end{codeblock}

% ------------------------------------------------------------
\section{Exercices}
\label{sec:ethique:exercices}
% ------------------------------------------------------------

\begin{exercice}[Analyse de biais]
T\'el\'echarger des embeddings Word2Vec ou FastText pr\'e-entra\^in\'es.
Calculer le score WEAT pour le biais de genre dans les professions.
Comparer les r\'esultats entre les embeddings entra\^in\'es sur des
corpus diff\'erents (Wikip\'edia vs Common Crawl).
\end{exercice}

\begin{exercice}[D\'ebiaisage]
Impl\'ementer la m\'ethode de d\'ebiaisage de Bolukbasi et al.
V\'erifier que les analogies biais\'ees sont r\'eduites. Tester si
le d\'ebiaisage affecte la performance sur une t\^ache de similarit\'e
s\'emantique (ex.~: STS-B).
\end{exercice}

\begin{exercice}[Toxicit\'e des LLM]
Utiliser l'API OpenAI Moderation (ou un classifieur Hugging Face) pour
\'evaluer la toxicit\'e de 100 g\'en\'erations d'un mod\`ele de langue
\`a partir d'amorces neutres. Quels types de toxicit\'e sont les plus
fr\'equents~?
\end{exercice}

\begin{exercice}[Model Card]
R\'ediger une model card compl\`ete pour un classifieur de sentiment
entra\^in\'e sur des critiques de films en fran\c{c}ais. Inclure les
limitations, les biais potentiels et les recommandations d'utilisation.
\end{exercice}
